% revised: CC 8 October 2026
% What this holds: Section Local Authority.
% Read by arlington-bsap.tex. Prose lives here; formatting lives in the wrapper and style/.
\section{Local Authority}\label{sec:legal}

The sections before this one read the Board's history through the equity
resolution, asking who has held its seats and how it divides its work with the
Manager. One matter sits outside that lens: what legal authority Virginia leaves
Arlington to decide for itself. By agreement with County Board staff, this report
treats it briefly. The County Attorney's July 2026 report is the primary
resource, and this report adds only the history.\autocite{samuel2026}

Virginia law, not the Board's own choice, sets what Arlington may decide for
itself. The state's courts took that rule from Dillon on Municipal Corporations: a
local government holds only the powers its charter grants in words or by
necessary implication. \textit{Kirkham v. Russell} (1882) applied it to
Petersburg's council, and \textit{City of Winchester v. Redmond} (1896) quoted
Dillon outright, which together made Dillon's Rule Virginia
law.\autocite{kirkhamvrussell1882,winchestervredmond1896}\fnsep\autocite[195--96]{schragger2022}

% The 1958-59 Commission on Incorporation is out of the report until a
% source says what it recommended: waits on: incorporation-commission-1959

Home rule was tried and failed. In the 1968--71 revision of the state
constitution, the revision commission proposed letting a charter county or city
act on anything state law did not deny it, and the General Assembly struck the
provision before the new constitution reached voters.\autocite[200]{schragger2022}

The County Attorney's report lists three decisions under the rule, and all three
are Arlington's own. In 1977 the Virginia Supreme Court voided the County's
agreements with its employees and barred public-sector bargaining statewide, a
bar the General Assembly lifted only in
2021.\autocite{commonwealthvarlington1977} In 2000 it held that the County had no
authority to extend health benefits to employees' domestic
partners.\autocite{arlingtonvwhite2000} The 1985 decision is the one that bears on
duties. The Board accepted a letter of intent to lease 6.4 acres of the courthouse
parking lot to a developer for 75 years and directed its Manager to sign. He
refused, and the Board sued to compel him. The Virginia Supreme Court denied the
writ because a county has no power to lease, so the Board had never held the
power it asked its own Manager to exercise.\autocite{brownvarlington1985} Part C's
Board Duties section returns to \textit{County Board of Arlington v. Brown}.

